\begin{figure}[tbp]
\centering
\begin{figurealt}{Porting begins by freezing the original game's behavioral and asset baseline. Translate API shape mechanically and compile against CNA's strict XNA surface. Inventory every asset and select a concrete XNB, CNJ, native, or glTF route. Choose renderer and shader strategies, adapt ownership and services, then compare behavior and pixels against the baseline. A mismatch loops back to the smallest responsible boundary before deployment across the intended platform matrix.}
\resizebox{0.96\textwidth}{!}{%
\begin{tikzpicture}
  \node[cna evidence, text width=34mm] (baseline) at (-5.4,2.2) {freeze original baseline};
  \node[cna internal, text width=30mm] (translate) at (-1.7,2.2) {mechanical API\\translation};
  \node[cna public, text width=28mm] (strict) at (1.7,2.2) {strict-XNA compile};
  \node[cna internal, text width=30mm] (assets) at (5.0,2.2) {asset-route manifest};
  \node[cna internal, text width=30mm] (renderer) at (5.0,0) {renderer + shader\\strategy};
  \node[cna internal, text width=36mm] (adapt) at (1.5,0) {services, ownership,\\lifetime adaptation};
  \node[cna evidence, text width=31mm] (compare) at (-2.2,0) {behavior + pixel\\comparison};
  \node[cna public, text width=31mm] (deploy) at (-5.7,0) {target deployment\\matrix};
  \node[cna warning, text width=53mm] (localize) at (-2.2,-2.0) {localize mismatch; return to the\\API, asset, renderer, or service step};
  \draw[cna flow] (baseline) -- (translate);
  \draw[cna flow] (translate) -- (strict);
  \draw[cna flow] (strict) -- (assets);
  \draw[cna flow] (assets) -- (renderer);
  \draw[cna flow] (renderer) -- (adapt);
  \draw[cna flow] (adapt) -- (compare);
  \draw[cna flow] (compare) -- (deploy);
  \draw[cna optional] (compare) -- node[cna label]{mismatch} (localize);
\end{tikzpicture}}
\end{figurealt}
\caption{A porting workflow that preserves attribution. Translation, asset conversion, renderer
choice, and service adaptation are verified as separable changes.}
\label{fig:porting-workflow}
\end{figure}
